\documentclass[12pt]{article}

%%%%%%%%%%%%%%%%%%%%%%%%%%%%%%%%%%%%%%%%%%%%%%%%%%%%%%%%%%%%%%%%%%
%                     PREAMBLE / PACKAGES
%%%%%%%%%%%%%%%%%%%%%%%%%%%%%%%%%%%%%%%%%%%%%%%%%%%%%%%%%%%%%%%%%%

% ----------------------------------------------------------------
% Font encoding and page geometry
% ----------------------------------------------------------------
\usepackage[T1]{fontenc}
\usepackage[margin=1in]{geometry}

% ----------------------------------------------------------------
% Line and paragraph spacing
% ----------------------------------------------------------------
\usepackage{setspace}
\usepackage{parskip}

% ----------------------------------------------------------------
% Section and table-of-contents formatting
% ----------------------------------------------------------------
\usepackage{titlesec}
\usepackage{tocloft}

% ----------------------------------------------------------------
% Graphics, figures, and page layouts
% ----------------------------------------------------------------
\usepackage{graphicx}
\usepackage[font=small]{caption}
\usepackage{subcaption}
\usepackage{float}
\usepackage{placeins}

% ----------------------------------------------------------------
% Mathematics
% ----------------------------------------------------------------
\usepackage{amsmath}
\usepackage{amssymb}
\usepackage{amsthm}

% ----------------------------------------------------------------
% Tables and lists
% ----------------------------------------------------------------
\usepackage{booktabs}
\usepackage{tabularx}
\usepackage{enumitem}

% ----------------------------------------------------------------
% Citations and hyperlinks
% ----------------------------------------------------------------
\usepackage{hyperref}

% ----------------------------------------------------------------
% Default line spacing
% ----------------------------------------------------------------
\onehalfspacing

\title{Lab 11}
\author{Chuan Wei Heng}
\date{\today}

\begin{document}

\maketitle

\section*{2. Before Lab}

We begin by implementing the composite trapezoidal and
Simpson's rules. Both approximate an integral by dividing
an interval into smaller pieces, but Simpson's rule uses
quadratic approximations rather than straight lines.

For $M$ equally spaced points, the spacing is
$h=(b-a)/(M-1)$. The composite trapezoidal rule is

\[
I\approx h\left[\frac{f(a)+f(b)}{2}
+\sum_{j=1}^{M-2}f(x_j)\right].
\]

For Simpson's rule, we use an even number of subintervals.
In this lab, $M=5$, so there are four subintervals
within each application of the rule. We also compare
these methods with the supplied Gauss--Legendre quadrature.

\section*{3. Lab Day}

Next, we apply adaptive quadrature. The idea is to
approximate the integral on an interval, split that interval
in half, and compare the result with the sum of the two
smaller integrals. The difference gives an estimate of
the error:

\[
E=\left|I_{\mathrm{left}}+I_{\mathrm{right}}
-I_{\mathrm{whole}}\right|.
\]

If this is larger than the chosen tolerance, the interval
is divided again. Otherwise, we accept the result and
move on. We use the same adaptive algorithm for all three
quadrature rules, with five nodes per subinterval and
a tolerance of $10^{-3}$.

\section*{4. Exercises}

We now approximate

\[
I=\int_{0.1}^{2}\sin\left(\frac{1}{x}\right)\,dx.
\]

A reference value obtained using SciPy is approximately
$I=1.1455808341$. The results from the three methods were

\[
\begin{aligned}
I_{\mathrm{trap}} &\approx 1.1453764,\\
I_{\mathrm{Simpson}} &\approx 1.1455955,\\
I_{\mathrm{Gauss}} &\approx 1.1455814.
\end{aligned}
\]

The corresponding absolute errors were approximately
$2.04\times10^{-4}$, $1.47\times10^{-5}$, and
$5.51\times10^{-7}$, respectively. All three methods
therefore achieved the requested accuracy.

\begin{figure}[H]
    \centering
    \includegraphics[width=0.85\linewidth]{figure/adaptive_intervals.png}
    \caption{Adaptive intervals used by the three quadrature methods.}
    \label{fig:adaptive_intervals}
\end{figure}

From the figure, we can see that the methods use more
intervals near $x=0.1$, where the function oscillates
more rapidly. The composite trapezoidal rule required
9 intervals, Simpson's rule required 6, and Gauss--Legendre
required only 5.

\begin{figure}[H]
    \centering
    \includegraphics[width=0.85\linewidth]{figure/absolute_errors.png}
    \caption{Absolute errors for the three adaptive quadrature methods.
    The dashed line indicates the required tolerance.}
    \label{fig:quadrature_errors}
\end{figure}

Interestingly, Gauss--Legendre quadrature performs best
for this example, both in terms of the number of intervals
used and the resulting accuracy. This makes sense since
Gaussian quadrature can integrate higher-degree polynomials
exactly with the same number of nodes.

\section*{Conclusion}

Overall, we implemented three adaptive quadrature methods
and compared their performance on the same integral.
Gauss--Legendre was the best choice in this experiment,
requiring only 5 intervals to achieve the desired accuracy,
compared with 6 for Simpson's rule and 9 for the
trapezoidal rule.

This also demonstrates why adaptive quadrature is useful:
we can concentrate more points where the function changes
rapidly, instead of using a uniformly fine grid everywhere.

For details of the code, we turn the reader towards the
Git repository by \textit{chuanweiheng-cell} entitled
APPM4600. Thank you!

\end{document}
